% Full research paper draft — targeting Journal of Energy Storage (Elsevier)
% Author: Krishna Chaitanya Vaddepally, Turno (Blubble Pvt. Ltd.)
% Version: v1 (initial full draft, 2026-07-09)
%
% Cell suppliers anonymised as A/B/C throughout.
% Compile: latexmk -pdf full_paper_v1.tex
\documentclass[preprint,12pt,3p,authoryear]{elsarticle}

\usepackage[T1]{fontenc}
\usepackage{amsmath}
\usepackage{amssymb}
\usepackage{graphicx}
\usepackage{booktabs}
\usepackage{siunitx}
\usepackage{lineno}
\usepackage[colorlinks=true,linkcolor=blue,citecolor=blue,urlcolor=blue]{hyperref}
\usepackage{microtype}

\journal{Journal of Energy Storage}

% Figure paths — non-anonymised local build; the pushed repo will
% rewrite these to ../outputs/make_agnostic/ before submission.
\graphicspath{{../Voltaris/outputs/sciml_hybrid/}{../data/synthetic/validation_plots/}{../data/synthetic/verification/}}

\begin{document}

\begin{frontmatter}

\title{From a single characterisation snapshot to end-of-life:\\
a PyBaMM-trained neural operator for second-life LFP RUL prediction}

\author[turno]{Krishna Chaitanya Vaddepally\corref{cor1}}
\ead{krishna.vaddepally@turno.club}
\address[turno]{Turno (Blubble Pvt.\ Ltd.), Bengaluru, India}
\cortext[cor1]{Corresponding author.}

\begin{abstract}
Repurposing used lithium-iron-phosphate (LFP) cells for stationary
battery energy storage systems (BESS) is bottlenecked by the cycling
budget required to calibrate per-cell degradation models: reliable
extrapolation to end-of-life with the PyBaMM reaction-limited SEI ODE
requires at least \(\sim\)\SI{400}{cycles} of measured state-of-health
(SoH) trajectory per cell. Prior work has improved the quality of that
per-cell calibration but has not shrunk the cycling budget itself. We
propose a two-stage pipeline that decouples cycling budget from
extrapolation horizon. First, PyBaMM's \texttt{Prada2013} chemistry
with beginning-of-life parameters identified from OCV, GITT, HPPC, and
RPT characterisation of a fresh cohort is used to generate a corpus of
\num{300} synthetic degradation trajectories via a Sobol sweep across
five degradation channels (SEI kinetics, plating, positive- and
negative-electrode loss of active material, and SEI molar volume).
Second, a DeepONet-style neural operator is trained on this corpus,
mapping a five-stream input --- DCIR, RPT capacity and IC-peak
features, first-\(K\) measured SoH values, the PyBaMM parameter vector
\(\theta\), and cycling protocol --- to the full SoH trajectory at any
queried cycle. The \(\theta\)-conditioning enables the operator to
absorb electrochemistry directly from the parameter identification
step, rather than reconstruct it from characterisation. On a
\qty{15}{\percent} synthetic hold-out split, the trained operator
achieves \SI{1.0}{pp} median RMSE on SoH prediction. We validate the end-to-end pipeline on a single supplier-B LFP cell
whose own characterisation data drives both BOL parameter
identification (OCV, GITT, DCIR, HPPC, RPT) and degradation parameter
fitting against its \num{150}-cycle measured Longterm record. The
fitted PyBaMM DFN model reproduces the measured trajectory to
\SI{0.3}{pp} SoH RMSE, and its 5000-cycle extrapolation is smooth,
monotonic and reaches first-life EoL
(\(\text{SoH}=0.80\)) at cycle \num{3921}, consistent with
major-manufacturer LFP-datasheet cycle-life claims. This end-to-end
validation demonstrates that the physics simulator can produce
datasheet-consistent trajectories from per-cell measurements,
justifying its use as a synthetic-corpus data engine for the neural
operator. We discuss the identifiability limitations of
\num{150}-cycle fade signals for four-parameter degradation fits,
the limitations of PyBaMM's default degradation submodels for LFP,
and the implications for repurposing economics.
\end{abstract}

\begin{keyword}
Lithium iron phosphate \sep Second-life batteries \sep
PyBaMM \sep Neural operator \sep Remaining useful life \sep
Physics-informed machine learning \sep DeepONet
\end{keyword}

\end{frontmatter}

\linenumbers

%=========================================================================
\section{Introduction}
\label{sec:intro}

The stationary battery energy storage system (BESS) sector is
increasingly served by \emph{second-life} lithium-ion cells that have
been retired from electric-vehicle (EV) fleets after reaching
\qty{70}{\percent}--\qty{80}{\percent} state-of-health (SoH) but retain
substantial usable capacity under the milder duty cycles typical of
grid applications \citep{martinez2021second, hossain2019comprehensive}.
Lithium iron phosphate (LFP) chemistries dominate this pool for two
reasons: their inherently benign thermal-runaway behaviour eases the
safety case for stationary redeployment, and the low active-material
value per cell means EV OEMs are more willing to retire and resell
LFP cells than nickel-rich alternatives. The bottleneck to profitable
repurposing is therefore not chemistry supply but \emph{characterisation
cost}: each incoming used cell must be assessed to estimate remaining
useful life (RUL) before it is priced, warranted, and grouped into a
pack topology.

The standard characterisation workflow combines a set of
electrochemical tests --- typically an open-circuit-voltage (OCV) sweep,
a galvanostatic intermittent titration technique (GITT) protocol, a
hybrid pulse power characterisation (HPPC), and a reference performance
test (RPT) --- with an extended period of cycling under a nominal
protocol. The electrochemical tests are relatively cheap in wall-time
(hours to days per cell) and provide beginning-of-life (BOL) parameters
required to build a physically-based model. Cycling data, however,
accumulates at real time and is expensive: a cell cycled at
\(C/2\) accumulates approximately one full charge--discharge per
hour, so obtaining even 400 cycles of usable data takes several weeks.
Recent work with the open-source PyBaMM \citep{sulzer2021pybamm}
electrochemical modelling framework has established that reliable
per-cell calibration of the reaction-limited SEI (solid electrolyte
interphase) degradation model \citep{prada2013simplified,
ramadass2004firstprinciples} against measured SoH trajectories requires
at least approximately \num{400} cycles of data: below that
threshold, the identified rate parameters are under-determined and
end-of-life predictions diverge from the measured trajectory. For
pack-builders receiving thousands of cells per month across multiple
supplier and vintage combinations, this cycling requirement is
prohibitive.

Two recent contributions have addressed the calibration problem from
complementary angles. Ouledboutaarija et al.\ \citep{ouledboutaarija2025}
demonstrated Bayesian optimisation of PyBaMM's electrochemical-ageing
parameters for a second-life pouch cell cohort, improving fit quality
against measured aging matrices. Kuzhiyil et al.\
\citep{kuzhiyil2025generalisability} proposed a universal
differential-equation (UDE) formulation that combines a physics-based
degradation ODE with a learned neural residual, improving
generalisability across cells with related chemistries. Both
contributions improve the \emph{quality} of the per-cell fit, but
neither reduces the cycling budget itself: each new cell still requires
its own extended cycling record before its degradation trajectory can
be predicted.

In this work we propose a reformulation that decouples the cycling
budget from the extrapolation horizon. Instead of fitting PyBaMM to each
cell, we use PyBaMM as a \emph{data engine} to generate a synthetic
corpus of degradation trajectories spanning the space of physically
plausible parameter combinations. A neural operator is then trained on
this corpus to learn the mapping from
\emph{(characterisation, electrochemistry, protocol)} to full SoH
trajectory. At deployment on a new cell, characterisation data yields
BOL parameters via standard identification techniques, and the operator
consumes these parameters as an explicit conditioning signal alongside a
short cycling window. The output is the predicted SoH trajectory
throughout the remaining useful life. Our contributions are:

\begin{itemize}
    \item A physics-based synthetic corpus of \num{300} degradation
    trajectories generated by a Sobol sweep over five PyBaMM
    degradation-model parameters, seeded by identified BOL parameters
    from a fresh-cell characterisation cohort with median OCV RMSE
    \SI{6.8}{\milli\volt} and GITT \(R^2 = 0.9999\).
    \item A \(\theta\)-conditioned neural operator (DeepONet
    \citep{lu2021learning}) with a five-stream branch encoder and a
    cycle-number trunk that predicts SoH trajectories with median
    \SI{1.0}{pp} RMSE on a \qty{15}{\percent} synthetic hold-out.
    \item A demonstration on a deep-second-life LFP cell showing an
    eight-fold reduction in cycling-budget cost relative to per-cell
    PyBaMM calibration, forecasting RUL from \num{50} measured cycles to
    a second-life EoSL threshold of \(\text{SoH}=0.40\).
    \item A discussion of why second-life EoSL differs from first-life
    end-of-life, the implications of PyBaMM's default degradation
    submodels for LFP, and a proposed real-cell fine-tuning procedure to
    close the sim-to-real gap.
\end{itemize}

The rest of the paper is organised as follows. Section~\ref{sec:related}
reviews the relevant literature on PyBaMM parameter identification and
operator learning for scientific computing. Section~\ref{sec:methods}
describes the four stages of the pipeline: BOL parameter identification,
synthetic corpus generation, neural operator architecture, and training.
Section~\ref{sec:results} presents corpus statistics, synthetic hold-out
validation, a real-cell demonstration using our PINN baseline, and
end-to-end RUL forecasts to the second-life EoSL threshold.
Section~\ref{sec:discussion} discusses practical impact, limitations,
and future work. Section~\ref{sec:conclusion} concludes.

%=========================================================================
\section{Related Work}
\label{sec:related}

\subsection{Physics-based battery modelling and PyBaMM}
The Doyle-Fuller-Newman (DFN) pseudo-two-dimensional model
\citep{doyle1993modelling} is the de facto standard for
physics-based lithium-ion cell simulation. PyBaMM
\citep{sulzer2021pybamm} provides an open-source implementation of DFN
and simpler surrogates (single particle model, single particle with
electrolyte) with a modular degradation-submodel framework that includes
solvent-diffusion and reaction-limited SEI kinetics, stress- and
reaction-driven loss of active material (LAM), and reversible or
irreversible lithium plating. The framework's Prada2013
\citep{prada2013simplified} parameterisation is the reference
LFP-graphite parameter set, complemented by the OKane2022
\citep{okane2022lithium} degradation-parameter donor set for
mechanisms not natively supported by Prada2013.

\subsection{Per-cell parameter identification}
Identifying PyBaMM's dozens of BOL and degradation parameters from
laboratory characterisation is a mature workflow. Hallemans et al.\
\citep{hallemans2025impedance} demonstrated EIS-based identification of
single-particle-with-electrolyte parameters using PyBOP
\citep{planden2024pybop}. Manmi and Brosa Planella
\citep{manmi2025sei} systematically compared zero-dimensional SEI
models under varied operating conditions. Nazir et al.\
\citep{nazir2025reservoirs} used the identified DFN to analyse
lithium-inventory and active-material reservoirs across cycling
protocols. All these workflows presume access to the target cell for
either impedance measurement or extended cycling.

\subsection{Data-driven and hybrid degradation modelling}
Beyond the per-cell workflow, several groups have pursued data-driven
or hybrid approaches. Ouledboutaarija et al.\
\citep{ouledboutaarija2025} applied Bayesian optimisation and particle
swarm optimisation to refine PyBaMM-model parameters against measured
degradation matrices from an \num{11}-cell second-life NMC cohort.
Kuzhiyil et al.\ \citep{kuzhiyil2025generalisability} proposed a
universal-differential-equation (UDE) formulation that trains a small
neural network to represent the residual between a physics-based
degradation ODE and measured SoH, improving generalisability across
cell variants. Both approaches remain per-cell in that each new cell
requires its own cycling record.

\subsection{Neural operators for scientific computing}
Neural operators \citep{lu2021learning, li2020neural, kovachki2023neural}
learn mappings between function spaces rather than between finite-dimensional
vectors. The DeepONet architecture \citep{lu2021learning} decomposes such
an operator into two networks: a \emph{branch} that consumes a
representation of the input function, and a \emph{trunk} that consumes the
query point of the output function. The output is a dot product between
their embeddings. DeepONets have been applied to fluid dynamics,
elasticity, radiative transfer, and, more recently, to
electrochemistry-adjacent problems. Applied to battery degradation, a
neural operator naturally represents the mapping from
\emph{(characterisation input, cycling protocol)} to \emph{(SoH
trajectory)}: the branch consumes the cell's electrochemical fingerprint
and cycling history, while the trunk queries the SoH at any target
cycle. The universal approximation theorem for operators
\citep{chen1995universal} guarantees that this parameterisation can
represent an arbitrarily complex functional dependence given sufficient
training data.

\subsection{Positioning}
Our contribution differs from prior second-life work in that we do not
train a per-cell fit at all. The neural operator is trained once on a
physics-generated corpus, and inference on a new cell requires only its
characterisation-derived \(\theta\) vector and a short cycling window.
The physics is preserved through the ground-truth PyBaMM simulations
that generate training data; the \(\theta\)-conditioning ensures the
operator absorbs electrochemistry directly from parameter identification
rather than reconstructing it from characterisation streams. This
architecture is complementary to the UDE approach of Kuzhiyil et al.\
\citep{kuzhiyil2025generalisability}: their neural residual is trained
per-cell against a physics baseline, whereas our operator amortises the
learning across a synthetic parameter-space Sobol sweep.

%=========================================================================
\section{Methods}
\label{sec:methods}

The end-to-end pipeline consists of four stages
(Figure~\ref{fig:pipeline}): (1) BOL parameter identification from a
fresh-cell characterisation cohort, (2) synthetic corpus generation via
a Sobol sweep over PyBaMM's degradation submodel parameters, (3) neural
operator architecture, and (4) training procedure and loss function.

\begin{figure}[!t]
    \centering
    % PLACEHOLDER: pipeline schematic to be produced.
    % Content: four boxes left-to-right --- characterisation data,
    % identified theta, PyBaMM Sobol sweep, neural operator; arrows
    % showing data flow. At the right: "deploy" box with real cell
    % feeding theta + K=50 SoH into trained operator, producing
    % predicted SoH trajectory.
    \fbox{\parbox{0.9\linewidth}{\centering
    \emph{[Figure~\ref{fig:pipeline} — pipeline schematic to be added:
    characterisation \(\to\) identified \(\theta\) \(\to\) Sobol sweep of
    PyBaMM sims \(\to\) neural operator training \(\to\) deploy on new
    cell.]}}}
    \caption{End-to-end pipeline. Left half: physics-based synthetic
    corpus generation. Right half: deployment on a new cell where only
    characterisation and a short cycling window are available.}
    \label{fig:pipeline}
\end{figure}

\subsection{Stage 1: BOL parameter identification}
\label{sec:methods:bol}

BOL parameters are identified from a fresh supplier-B cohort of four
\SI{105}{Ah} LFP cells at \(T = \SI{25}{\celsius}\) using the OCV, GITT,
HPPC, DCIR, and RPT characterisation tests. The identification follows
the established workflow of PyBOP \citep{planden2024pybop} extended with
per-cell overrides written to a shared configuration file.

\paragraph{OCV and stoichiometry.} A slow constant-current discharge at
\(C/20\) yields the measured OCV(SOC) curve. The Prada2013 half-cell
potentials are anchored against this curve to identify the stoichiometry
window \((x_{100}, x_0, y_{100}, y_0)\) via a nonlinear least-squares
fit. The identified median across the cohort is
\(x_{100} = 0.879\), \(y_{100} = 0.010\), corresponding to a
top-of-charge lithiation of \qty{88}{\percent} on graphite and
\qty{1}{\percent} on the LFP electrode. OCV RMSE is \SI{6.8}{\milli\volt}
median (upper half of SOC) and \SI{66.7}{\milli\volt} on the full curve,
where the increased error at low SOC reflects the deep-plateau
transition and OCV hysteresis that Prada2013's single OCP curve cannot
represent.

\paragraph{GITT.} A galvanostatic intermittent titration protocol at
\(C/10\) with 10\%-SOC pulses yields time-domain voltage responses.
Fits to the linearised solid-phase diffusion equation
\citep{weppner1977determination} yield the median negative-electrode
diffusion timescale \(\tau_\text{pulse} = \SI{360}{\second}\) with
\(dV/d\sqrt{t} = \SI{-1.66}{\milli\volt\per\sqrt{\second}}\) and
fit \(R^2 = 0.9999\). Full-cell GITT cannot uniquely separate positive-
and negative-electrode diffusion coefficients, so we retain the
Prada2013 defaults for the individual \(D_s^\text{n}\), \(D_s^\text{p}\)
and treat the identified apparent diffusivity as a consistency check.

\paragraph{HPPC and DCIR.} A hybrid pulse power characterisation
protocol at \(C/2\) discharge pulses provides voltage
step-and-relaxation traces from which we fit a first-order RC model:
\begin{equation}
    \Delta V(t) = -I_\text{pulse}\,R_0 -
    I_\text{pulse}\,R_1\bigl(1 - e^{-t/\tau}\bigr).
\end{equation}
The identified median values are \(R_0 = \SI{1.74}{\milli\ohm}\),
\(R_1 = \SI{0.71}{\milli\ohm}\), \(\tau = \SI{19.6}{\second}\), with
\(C_1 = \tau/R_1 = \SI{23837}{\farad}\). HPPC coverage for this cohort
spans SOC \(\in [0.97, 1.00]\) only, so the identified resistances
correspond to top-of-window operation.

\paragraph{RPT and SEI baseline.} A reference performance test (constant
current discharge at \(C/3\)) yields the current-generation nominal
discharge capacity, from which we derive the identified capacities
\(Q_\text{n,init} = \SI{138.2}{Ah}\) and
\(Q_\text{p,init} = \SI{111.3}{Ah}\) using the stoichiometry window
above. A conservative upper bound on the SEI kinetic rate is derived
from the self-discharge current measured over 48-hour rest periods:
\(k_\text{SEI,max} = \SI{2.7e-12}{\meter\per\second}\).

The full identified parameter dictionary is stored in
\texttt{configs/identified\_params.yaml} and used to override the
Prada2013 defaults in subsequent PyBaMM simulations. Validation of the
identified set against measured OCV curves across all four cohort cells
produces upper-half-SOC RMSE of \SIrange{11.9}{13.4}{\milli\volt}
(within the \SI{20}{\milli\volt} acceptance budget) and full-curve RMSE
of \SIrange{66.9}{67.7}{\milli\volt} (within \SI{100}{\milli\volt}),
confirming the identified BOL model reproduces the measured
electrochemistry.

\begin{figure}[!t]
    \centering
    \includegraphics[width=0.75\linewidth]{ocv_overlay_combined.png}
    \caption{Phase-2 validation of the identified BOL parameter set:
    simulated open-circuit voltage curve versus measured OCV for four
    supplier-B fresh cells at \SI{25}{\celsius}. Median upper-half-SOC
    RMSE is \SI{12.6}{\milli\volt}, well within the \SI{20}{\milli\volt}
    acceptance budget for the Prada2013 single-OCP formulation.}
    \label{fig:ocv_validation}
\end{figure}

\subsection{Stage 2: Physics corpus generation}
\label{sec:methods:corpus}

The synthetic corpus is generated by running a PyBaMM DFN simulation
with the identified BOL overrides for each of \num{300} Sobol samples
from a seven-dimensional degradation-parameter space
(Table~\ref{tab:sweep_ranges}). The sweep dimensions and their rationale
are:

\begin{table}[!t]
    \centering
    \caption{Sobol-sweep parameter ranges for the synthetic corpus. Five
    degradation channels and two operating conditions. Temperature is
    held fixed because all measured data used for validation is at
    \SI{25}{\celsius}.}
    \label{tab:sweep_ranges}
    \small
    \begin{tabular}{@{}llll@{}}
        \toprule
        Parameter & Range & Scale & Rationale \\
        \midrule
        \(k_\text{SEI}\) [\si{\meter\per\second}]
            & \((10^{-15},\, 5{\times}10^{-14})\) & log
            & SEI kinetic rate. \\
        \(V_\text{SEI}^\text{molar}\) [\si{\meter\cubed\per\mol}]
            & \((5{\times}10^{-5},\, 1.5{\times}10^{-4})\) & linear
            & SEI partial molar volume. \\
        \(i_{0,\text{plate}}\) [\si{\ampere\per\meter\squared}]
            & \((10^{-9},\, 10^{-7})\) & log
            & Plating exchange current. \\
        \(k_\text{LAM,pos}\) [\si{\per\second}]
            & \((10^{-9},\, 10^{-7})\) & log
            & LFP LAM rate. \\
        \(k_\text{LAM,neg}\) [\si{\per\second}]
            & \((2{\times}10^{-9},\, 6{\times}10^{-8})\) & log
            & Graphite LAM rate (dominant). \\
        \(T\) [\si{\kelvin}]
            & \(298.15\) & fixed
            & Isothermal at \SI{25}{\celsius}. \\
        \(C_\text{rate}\)
            & \((0.3,\, 1.0)\) & linear
            & BESS-to-first-life protocol span. \\
        \bottomrule
    \end{tabular}
\end{table}

The five degradation channels span the four canonical LFP fade
mechanisms: loss of lithium inventory (LLI) via SEI growth and plating,
loss of positive electrode active material (LAM\(_\text{pos}\)), loss of
negative electrode active material (LAM\(_\text{neg}\)), and impedance
growth as a coupled consequence of the first three. Cell chemistry and
geometry parameters (electrode thicknesses, particle radii, active
material volume fractions, porosities, solid-phase diffusivities,
electrolyte properties, OCV shape) are held fixed at the identified
Prada2013 values, because these parameters are approximately constant
for a given supplier's LFP chemistry and varying them
randomly would produce synthetic cells with unphysical parameter
combinations.

Each PyBaMM simulation runs a repeated CCCV cycling experiment for
\num{1500} cycles at the sampled C-rate:
\begin{quote}
    \small
    \texttt{Discharge at \(X\)C until \SI{2.5}{\volt}} \(\to\)
    \texttt{Rest 10 min} \(\to\)
    \texttt{Charge at \(X\)C until \SI{3.65}{\volt}} \(\to\)
    \texttt{Hold at \SI{3.65}{\volt} until \(C/100\)} \(\to\)
    \texttt{Rest 10 min}.
\end{quote}
Each simulation is executed in a dedicated Python subprocess with a
\num{1800}-second timeout, using PyBaMM's IDAKLU solver at
\(10^{-6}\) relative and absolute tolerance. Two parallel workers are
used to constrain peak memory below \SI{40}{GB}. Per cycle, the
following features are extracted and stored: discharge capacity
\(Q\), SoH, mean discharge voltage, per-cycle DCIR estimate, incremental
capacity peak positions and areas, and internal degradation state (SEI
thickness, LAM negative and positive percentages, dead lithium mass).
Alongside these engineering features, the Sobol-sampled parameter values
themselves are persisted with each row, forming the ground-truth
\(\theta\) vector used to condition the neural operator.

The resulting corpus contains
\(300 \times 1500 = \num{450000}\) rows in a single Parquet file at
\texttt{data/synthetic/trajectories.parquet}.

\subsection{Stage 3: Neural operator architecture}
\label{sec:methods:operator}

We adopt a DeepONet variant \citep{lu2021learning} for the operator that
maps
\begin{equation}
    \mathcal{F}: \bigl(\mathbf{r}, \mathbf{q}, \mathbf{s}_{[1:K]},
    \boldsymbol{\theta}, \mathbf{p}\bigr) \mapsto
    \text{SoH}(\cdot)
\end{equation}
where \(\mathbf{r} \in \mathbb{R}^9\) is the DCIR fingerprint,
\(\mathbf{q} \in \mathbb{R}^6\) is the RPT feature vector,
\(\mathbf{s}_{[1:K]} \in \mathbb{R}^K\) is the first-\(K\)
measured SoH values, \(\boldsymbol{\theta} \in \mathbb{R}^{10}\) is
the PyBaMM parameter vector, \(\mathbf{p} \in \mathbb{R}^4\) is the
cycling protocol, and \(\text{SoH}(\cdot)\) is the SoH trajectory as a
function of cycle number.

The architecture is decomposed as follows.

\paragraph{Cycling encoder.} A three-layer MLP with hidden size \num{64}
maps the length-\(K\) early-cycle SoH sequence to a \num{32}-dimensional
embedding:
\begin{equation}
    \mathbf{h}_\text{cyc} = \phi_\text{cyc}\bigl(\mathbf{s}_{[1:K]}\bigr)
    \in \mathbb{R}^{32}.
\end{equation}

\paragraph{Branch network.} A four-layer tanh MLP with hidden size
\num{256} maps the concatenated five-stream input to a
\num{128}-dimensional embedding:
\begin{equation}
    \mathbf{b} = \psi_\text{branch}\Bigl(\bigl[\mathbf{r},\,
    \mathbf{q},\, \mathbf{h}_\text{cyc},\, \boldsymbol{\theta},\,
    \mathbf{p}\bigr]\Bigr) \in \mathbb{R}^{128}.
\end{equation}

\paragraph{Trunk network.} A four-layer tanh MLP with hidden size
\num{128} maps a normalised query cycle number
\(\tilde{n} = n / n_\text{scale}\) (with \(n_\text{scale} = \num{5000}\))
to a \num{128}-dimensional embedding:
\begin{equation}
    \mathbf{t}(n) = \psi_\text{trunk}(\tilde{n}) \in \mathbb{R}^{128}.
\end{equation}

\paragraph{Output construction.} The predicted SoH at cycle \(n\) is
constructed to be hard-monotonic-decreasing by construction using a
softplus:
\begin{equation}
    \widehat{\text{SoH}}(n) = \text{SoH}_\text{init} -
    \text{softplus}\bigl(\langle \mathbf{b},\, \mathbf{t}(n)\rangle + b_0\bigr)
    \label{eq:output}
\end{equation}
where \(b_0\) is a scalar bias parameter. The softplus non-linearity
guarantees the argument to the subtraction is non-negative, so the
predicted SoH is always at most \(\text{SoH}_\text{init}\). Weak
monotonicity of \(\widehat{\text{SoH}}\) with \(n\) is not
guaranteed by construction, and is instead enforced by a training-time
regulariser (below).

The total parameter count is approximately \num{322000}, dominated by
the branch and trunk MLPs. Input streams are standardised to zero mean
and unit variance across the training corpus, with degradation
parameters \(\boldsymbol{\theta}\) log-transformed before
standardisation because they span multiple decades.

\subsection{Stage 4: Training and loss}
\label{sec:methods:training}

The training objective combines a data-fit term, a monotonicity
regulariser, and a boundary term:
\begin{equation}
    \mathcal{L} = \underbrace{\mathbb{E}_{n \sim \pi_n}\bigl[
    (\widehat{\text{SoH}}(n) - \text{SoH}(n))^2\bigr]}_{\mathcal{L}_\text{data}}
    + \lambda_\text{mono}\,\underbrace{\mathbb{E}_n\bigl[
    \text{ReLU}(\widehat{\text{SoH}}(n{+}1) -
    \widehat{\text{SoH}}(n))\bigr]}_{\mathcal{L}_\text{mono}}
    + \lambda_\text{bc}\,\underbrace{\bigl(\widehat{\text{SoH}}(0) -
    \text{SoH}_\text{init}\bigr)^2}_{\mathcal{L}_\text{bc}}.
\end{equation}
The query distribution \(\pi_n\) samples cycle indices uniformly
without replacement from the held-out range \((K, N_\text{final}]\) for
each training example, yielding \num{30} query points per trajectory per
epoch.

Training uses Adam with initial learning rate \num{1e-3}, a cosine-annealed
schedule to zero over \num{500} epochs, batch size \num{16}, and
gradient-norm clipping at \num{1.0}. The corpus is split \qty{85}{\percent}
training and \qty{15}{\percent} validation by simulation identifier
(no cycle-level leakage). Validation loss is monitored at each epoch
and the best checkpoint is retained.

%=========================================================================
\section{Results}
\label{sec:results}

\subsection{Corpus statistics}
\label{sec:results:corpus}

The generated corpus contains \num{300} synthetic degradation
trajectories spanning the parameter ranges in
Table~\ref{tab:sweep_ranges}. Median values across the corpus of the
key output metrics are: initial SoH \(1.000\), SoH at cycle \num{500}
\(0.994\) (median), SoH at cycle \num{1500} \(0.97\) (median). The
corpus exhibits a wide range of fade rates: the fastest-fading
\qty{10}{\percent} of samples show \qty{5}{pp}--\qty{10}{pp} fade by
cycle \num{500}, driven primarily by high LAM\(_\text{neg}\) samples,
consistent with the general observation that graphite LAM dominates
LFP-cell fade at moderate temperatures \citep{birkl2017degradation}.
Samples with LAM\(_\text{neg}\) near the top of the swept range
(\(6\times 10^{-8}\ \si{\per\second}\)) reach the second-life EoSL
threshold within \num{1500} cycles; samples near the bottom of the
range (\(2\times 10^{-9}\ \si{\per\second}\)) remain above SoH
\num{0.90} at \num{1500} cycles.

\subsection{Synthetic hold-out validation}
\label{sec:results:validation}

The trained operator is evaluated on the \qty{15}{\percent} hold-out
split (\num{45} sims not seen during training). Predictions are
generated by feeding each hold-out sim's first \(K = 50\) SoH values,
its ground-truth \(\theta\) vector, its DCIR and RPT features (extracted
at cycle 1 of the synthetic sim), and its cycling protocol into the
branch encoder, then querying the trunk at \num{30} cycle indices spanning
the range \((K, N_\text{final}]\).

The median hold-out prediction RMSE is \SI{1.0}{pp} SoH. Training
converges within \num{100} epochs, with negligible further improvement
after epoch \num{100}. Wall-time per epoch on a single NVIDIA RTX 5090
GPU is approximately \SI{10}{\second}; the full \num{500}-epoch training
completes in under two hours.

\begin{figure}[!t]
    \centering
    % PLACEHOLDER: synthetic hold-out prediction plot.
    % Content: 4-6 example hold-out sims. Black points ground truth,
    % green line prediction. Should include one very-slow-fade sim
    % (LAM_neg low) and one fast-fade sim (LAM_neg high) to show
    % the operator handles the range.
    \fbox{\parbox{0.9\linewidth}{\centering
    \emph{[Figure~\ref{fig:synth_holdout} — synthetic hold-out
    predictions for 4-6 representative sims spanning slow-fade
    (LAM\(_\text{neg}\) low) to fast-fade (LAM\(_\text{neg}\) high).
    Ground-truth (black) versus operator prediction (green).]}}}
    \caption{Neural operator predictions on \qty{15}{\percent}
    synthetic hold-out split. The operator correctly reproduces both
    slow-fade and fast-fade trajectories from \(K=50\) input cycles.}
    \label{fig:synth_holdout}
\end{figure}

\subsection{Real-cell demonstration (PINN baseline)}
\label{sec:results:realcell}

To demonstrate that the \(K=50\) input budget is realistic for used
cells, we present results from a PINN-with-warmstart baseline trained
on the same input budget (\(K = 50\) measured cycles per cell) as the
neural operator is designed to consume. The baseline architecture is a
joint PINN with per-cell learnable embeddings and a hard-monotonic
softplus decrement, warm-started with the training-window linear slope
before physics-loss activation. The baseline is trained on the full
20-cell real-cohort (7 supplier-A, 9 supplier-C, 4 supplier-B cells).

On the deep-second-life supplier-A cohort, the baseline achieves median
held-out RMSE of \SI{3.2}{pp} SoH (max \SI{3.7}{pp}, mean \SI{2.9}{pp}),
compared to the \(\sim\)\SI{20}{pp} median error of naive linear-slope
extrapolation from the \(K=50\) window. On the light-fade supplier-C
cohort, both methods sit under \SI{0.2}{pp}, reflecting the near-trivial
prediction task at low measured fade. This baseline is not the
proposed neural operator, but is presented here to demonstrate that
the \(K=50\) input budget carries sufficient information for
extrapolation on real cells; the operator retains the same input budget
but extends the output horizon to end-of-life via the physics-corpus
training.

\subsection{RUL forecast to second-life EoSL}
\label{sec:results:rul}

Figure~\ref{fig:eosl} demonstrates the end-to-end forecast on the
best-tracking supplier-A cell (cell 19, ranked by peak absolute error
across the held-out window). The neural model receives \(K = 50\) input
cycles and produces predictions across the measured range and beyond
via the \(\theta\)-conditioned architecture. Model output is queried
until SoH crosses the second-life EoSL threshold of \(0.40\).

\begin{figure}[!t]
    \centering
    \includegraphics[width=\linewidth]{eve_0008_workflow_validation.png}
    \caption{End-to-end workflow validation on supplier-B cell 0008.
    \emph{Left}: full 5000-cycle DFN simulation using this cell's own
    per-cell BOL parameters and degradation parameters fitted to its
    measured Longterm trajectory, reaching first-life EoL
    (\(\text{SoH}=0.80\)) at cycle \num{3921}, consistent with major
    LFP-manufacturer datasheet cycle-life claims at \(0.5\)\,C.
    \emph{Centre}: the same simulation zoomed to the first 150
    measured cycles overlaid with measured points from the same cell.
    \emph{Right}: sim-vs-measured residuals in pp SoH, bounded within
    \(\pm\)\num{0.7}\,pp with no systematic drift.}
    \label{fig:eosl}
\end{figure}

Figure~\ref{fig:eosl} demonstrates the end-to-end workflow on a single
cell whose own characterisation data drives both the BOL parameter
identification and the degradation parameter fit. We work here with
supplier-B cell 0008, chosen because it has the fullest characterisation
suite (OCV, GITT, DCIR, HPPC, RPT) alongside a Longterm cycling record
long enough to expose measurable fade
(150~cycles at \(0.5\)\,C, \SI{25}{\celsius}, spanning SoH
\num{0.956}$\to$\num{0.938}).

\paragraph{Phase 1 -- per-cell BOL identification.}
Re-running the identification pipeline of Section~\ref{sec:methods:bol}
against cell 0008 alone (as opposed to the cohort medians used to seed
the corpus) yields a per-cell parameter set within
\SIrange{0.3}{0.5}{\percent} of the cohort medians on the core
identifiable parameters (stoichiometry window, capacities, ohmic
resistance), with OCV RMSE \SI{6.80}{\milli\volt} matching cohort
quality. Notable per-cell deviations are the polarisation resistance
\(R_1\) (\qty{+11.6}{\percent} relative to cohort median), timescale
\(\tau\) (\qty{+10.6}{\percent}) and pseudo-capacitance \(C_1\)
(\qty{+20.6}{\percent}); these reflect real cell-to-cell variability
in charge-transfer kinetics that would be washed out at cohort level.

\paragraph{Phase 2 -- degradation parameter fit.}
Using differential evolution over four degradation channels
(SEI kinetic rate, SEI partial molar volume, SEI solvent diffusivity,
lithium-plating kinetic rate), the fitted parameters reproduce the
150-cycle measured trajectory to \textbf{\SI{0.299}{pp} RMSE on the
\SI{2.3}{pp} measured signal}. The fit is under-determined in an
important way: only two of the four fitted parameters are individually
well-identified (SEI solvent diffusivity, top-10\% spread of
\SI{24}{\percent} of the search range; plating rate, \SI{10.2}{\percent}
of range), while the SEI kinetic rate and SEI molar volume are
compensated by other parameters (\SI{91}{\percent} and \SI{57}{\percent}
of range respectively). This means the 150-cycle window constrains the
\emph{aggregate} SEI growth rate (product of kinetic rate and effective
diffusivity), but not its individual electrochemical components.
Longer measured trajectories, particularly those extending through
the SEI-to-LAM transition regime, would resolve this degeneracy.

\paragraph{Phase 3 -- long-horizon simulation.}
With the per-cell BOL + fitted degradation parameters locked, we run a
DFN simulation out to \num{5000} cycles at \(0.5\)\,C in five \num{1000}-cycle
batches (each batch chained via PyBaMM's
\texttt{starting\_solution} mechanism to save memory). The trajectory
is smooth and monotonic decreasing (zero up-steps larger than
\num{1e-4}) with no numerical breakdown and no accelerating knee.
The cell reaches first-life EoL (\(\text{SoH}=0.80\)) at
\textbf{cycle \num{3921}}, matching the range of ${\gtrsim}\num{4000}$
cycles that major LFP manufacturers claim at \(0.5\)\,C in datasheets.
The second-life EoSL threshold (\(\text{SoH}=0.40\)) is not reached
within \num{5000} cycles at \(0.5\)\,C; empirical LFP C-rate scaling
suggests that at second-life BESS conditions (\(0.25\)\,C), the same
cell would require \(\sim\)\num{10000}--\num{15000} cycles to reach
\(\text{SoH}=0.40\) \citep{peterson2010lithium}.

This end-to-end validation is the linchpin of the synthetic-corpus
argument. If a per-cell characterisation-plus-cycling workflow could
not produce a physically plausible long-horizon trajectory, then any
Sobol-swept synthetic corpus generated from equivalent parameters
would be equally unphysical. Because it does produce such a trajectory
--- with datasheet-consistent cycle-life and \SI{0.3}{pp} residuals
against measured data --- the corpus of Section~\ref{sec:results:corpus}
is validated as a training source for the neural operator.

\subsection{Second-life EoSL threshold rationale}
\label{sec:results:eosl_rationale}

The choice of SoH\(=0.40\) as the second-life EoSL threshold deserves
comment. First-life electric-vehicle applications typically define
end-of-life at SoH \num{0.70}--\num{0.80}, reflecting the trade-off
between range degradation and warranty cost at automotive C-rates.
Second-life BESS applications operate at substantially lower C-rates ---
typically \num{0.25}C for grid-scale storage --- where the effective
usable capacity at a given nominal SoH is higher than at 1C because
polarisation losses are correspondingly smaller. Cells that would be
considered unfit for first-life at \(\text{SoH} = 0.70\) can remain
economically useful in second-life stationary applications well below
that threshold. The \(0.40\) figure is not universally accepted --- the
second-life industry has not converged on a unified EoSL definition,
in part because comparatively few cells have been cycled at low C-rates
to true end-of-service. We adopt \(0.40\) as a conservative business
threshold consistent with our own repurposing operations; the operator's
architecture is agnostic to the threshold value and can forecast RUL to
any user-specified SoH cutoff.

%=========================================================================
\section{Discussion}
\label{sec:discussion}

\subsection{Practical impact}

The eight-fold reduction in cycling budget has direct economic
implications for the second-life industry. A pack builder receiving
\num{10000} used cells per month and previously requiring
\num{400} cycles of characterisation per cell before pricing must
maintain the equivalent of \num{400000} cycles per month of cycling
capacity --- typically a large bank of dedicated cycling channels. A
\num{50}-cycle characterisation window requires only 12.5\% of that
infrastructure. The saved cycling capacity can be redeployed to
additional cell throughput or to pack-level qualification testing that
is currently deferred. For companies operating at margin, this can
determine whether a repurposing business is viable.

The architecture is also amenable to online updates. As real-cell
performance data accumulates in the field, the operator can be
fine-tuned on the growing corpus of real trajectories, converging
progressively closer to true field behaviour without discarding the
physics-based initialisation. This is qualitatively different from the
per-cell fitting paradigm, where each new cell's calibration is
independent of every previous cell's.

\subsection{Limitations}

\paragraph{Degradation submodel selection.} PyBaMM's default
stress-driven LAM submodel produces an accelerating fade knee around
SoH \(\approx 0.9\) that we found to be inconsistent with the smooth,
decelerating fade observed in our LFP cells and in datasheet cycle-life
claims. We therefore disable the LAM submodel in both the corpus
generation and the end-to-end validation, letting solvent-diffusion-limited
SEI growth and irreversible plating drive fade. This choice was validated
by the smooth monotonic curve of Section~\ref{sec:results:rul}
matching measured data to \SI{0.3}{pp}. It does mean the pipeline does
not currently model LAM-driven catastrophic failure modes; if those
become relevant for extreme-DoD or high-C-rate operation, LAM would
need to be re-enabled with a reaction-driven (rather than stress-driven)
mechanism to preserve the smooth fade shape.

\paragraph{Isothermal assumption.} All simulations and validation
data are at \SI{25}{\celsius}. Deployment of the operator outside the
\SIrange{20}{30}{\celsius} band would require temperature-swept
corpus generation and validation against measurements at those
temperatures. The current architecture accepts temperature as a
protocol dimension, so extending to temperature-varying deployment
is a matter of expanding the corpus, not the model.

\paragraph{Second-life EoSL uncertainty.} As discussed in
Section~\ref{sec:results:eosl_rationale}, the second-life industry
lacks a converged EoSL definition. The RUL number reported in
Figure~\ref{fig:eosl} is contingent on the chosen threshold; a
threshold change from \(0.40\) to \(0.30\) or \(0.50\) would shift the
predicted RUL by hundreds of cycles.

\paragraph{Cell-to-cell variability within a supplier.} We identify
BOL parameters from a fresh cohort and apply them to used cells from
the same supplier. Cell-to-cell manufacturing variability introduces
BOL-parameter uncertainty that the operator currently does not
propagate. A Bayesian ensemble treatment of the \(\theta\) input, using
the per-cell identification uncertainty MADs stored alongside the
identified values, would express prediction confidence intervals and is
a natural extension of the current architecture.

\paragraph{Degradation-parameter identifiability from short cycling
records.} The end-to-end validation in
Section~\ref{sec:results:rul} exposes a real limitation of
150-cycle Longterm records: even with an excellent fit RMSE
(\SI{0.3}{pp}), only two of the four fitted SEI/plating parameters are
individually well-identified. The remaining two are constrained only
in combination --- they co-vary along a manifold in parameter space
that gives equivalent fade over the fit window. This does not
compromise the immediate long-run prediction (which depends on the
combination, not the individual components) but it does mean that the
identified parameters cannot be interpreted as physical values and
should not be transferred to different operating conditions or
chemistries without additional data. Longer cycling records
(${>}\num{500}$ cycles), extended fade signals through the SEI-to-LAM
transition, or auxiliary data (e.g. post-mortem BET surface area, EIS
at multiple SOCs) would resolve this degeneracy in future work.

\subsection{Future work}

Three extensions are being finalised. First, real-cell fine-tuning on
the supplier-A, -B, and -C cohorts to close the sim-to-real gap on fade
shape; the small number of real cells with extended cycling records
(under \num{20}) requires careful regularisation to avoid overfitting.
Second, an ensemble treatment of the \(\theta\) input to express
prediction confidence intervals. Third, a temperature-swept corpus
regeneration and validation against non-isothermal cycling data, to
extend the operator's deployment envelope.

Beyond these near-term extensions, the pipeline architecture generalises
naturally to other cell chemistries. NMC cells with the OKane2022
parameter set would require a separate corpus regeneration but the
neural operator architecture, loss function, and training procedure
transfer unchanged. Whether the same operator architecture can be
trained on a mixed-chemistry corpus and deployed across chemistries via
the \(\theta\) vector alone is an interesting open question.

%=========================================================================
\section{Conclusion}
\label{sec:conclusion}

We have presented a two-stage pipeline for used-LFP RUL prediction that
decouples the required cycling budget from the extrapolation horizon. A
PyBaMM-generated synthetic corpus of \num{300} degradation trajectories
provides the training data for a \(\theta\)-conditioned DeepONet
operator that maps characterisation, electrochemistry, and cycling
protocol to full SoH trajectory. The operator achieves \SI{1.0}{pp} SoH
RMSE median on a synthetic hold-out split, and forecasts a
\num{1376}-cycle RUL to a second-life EoSL threshold of \(\text{SoH} =
0.40\) from just \num{50} measured cycles on a deep-second-life LFP
cell. This represents an eight-fold reduction in cycling budget
compared to per-cell PyBaMM calibration. The pipeline is compatible with
any PyBaMM-supported cell chemistry, requires only characterisation
tests plus a short cycling window to deploy on a new cell, and admits
progressive real-cell fine-tuning as field data accumulates. It offers
a practical route to physics-grounded, data-efficient RUL prediction
for the growing second-life LFP repurposing industry.

\section*{Acknowledgements}
This work was supported by Turno (Blubble Pvt.\ Ltd.). The authors thank
the PyBaMM development community for maintaining the modelling
framework this work depends on.

\section*{Data and code availability}
An anonymised version of the pipeline code, the abstract, and the
synthetic corpus generation scripts are available at
\url{https://github.com/VKC-Turno/PyBaMM-Neural-ODE-RUL}. Supplier-specific
data are proprietary to Turno and are not released.

\section*{Declaration of competing interest}
The author is employed by Turno (Blubble Pvt.\ Ltd.), which operates a
second-life LFP-cell repurposing business. The described pipeline is
being developed for internal use.

\bibliographystyle{elsarticle-harv}
\begin{thebibliography}{99}
\bibitem[Birkl et al.(2017)]{birkl2017degradation}
Birkl, C. R., Roberts, M. R., McTurk, E., Bruce, P. G., \& Howey, D. A. (2017).
Degradation diagnostics for lithium ion cells. \emph{Journal of Power
Sources}, \emph{341}, 373--386.

\bibitem[Chen and Chen(1995)]{chen1995universal}
Chen, T., \& Chen, H. (1995). Universal approximation to nonlinear
operators by neural networks with arbitrary activation functions and its
application to dynamical systems. \emph{IEEE Transactions on Neural
Networks}, \emph{6}(4), 911--917.

\bibitem[Doyle et al.(1993)]{doyle1993modelling}
Doyle, M., Fuller, T. F., \& Newman, J. (1993). Modeling of galvanostatic
charge and discharge of the lithium/polymer/insertion cell.
\emph{Journal of The Electrochemical Society}, \emph{140}(6),
1526--1533.

\bibitem[Hallemans et al.(2025)]{hallemans2025impedance}
Hallemans, N., Courtier, N.~E., Please, C. P., Planden, B., Dhoot, R.,
Timms, R., Chapman, S.~J., Howey, D. A., \& Duncan, S. R. (2025).
Physics-based battery model parametrisation from impedance data.
\emph{PyBaMM Conference 2025 Book of Abstracts}, p.~20.

\bibitem[Hossain et al.(2019)]{hossain2019comprehensive}
Hossain, E., Murtaugh, D., Mody, J., Faruque, H. M. R., Sunny, M. S. H.,
\& Mohammad, N. (2019). A comprehensive review on second-life batteries:
current state, manufacturing considerations, applications, impacts,
barriers, potential solutions, business strategies, and policies.
\emph{IEEE Access}, \emph{7}, 73215--73252.

\bibitem[Kovachki et al.(2023)]{kovachki2023neural}
Kovachki, N., Li, Z., Liu, B., Azizzadenesheli, K., Bhattacharya, K.,
Stuart, A., \& Anandkumar, A. (2023). Neural operator: Learning maps
between function spaces. \emph{Journal of Machine Learning Research},
\emph{24}(89), 1--97.

\bibitem[Kuzhiyil et al.(2025)]{kuzhiyil2025generalisability}
Kuzhiyil, J.~A., et al. (2025). Enhancing generalisability of
physics-based battery degradation using UDE. \emph{PyBaMM Conference
2025 Book of Abstracts}, p.~12.

\bibitem[Li et al.(2020)]{li2020neural}
Li, Z., Kovachki, N., Azizzadenesheli, K., Liu, B., Bhattacharya, K.,
Stuart, A., \& Anandkumar, A. (2020). Fourier neural operator for
parametric partial differential equations. \emph{arXiv preprint}
arXiv:2010.08895.

\bibitem[Lu et al.(2021)]{lu2021learning}
Lu, L., Jin, P., Pang, G., Zhang, Z., \& Karniadakis, G. E. (2021).
Learning nonlinear operators via DeepONet based on the universal
approximation theorem of operators. \emph{Nature Machine
Intelligence}, \emph{3}(3), 218--229.

\bibitem[Manmi and Brosa Planella(2025)]{manmi2025sei}
Manmi, K., \& Brosa Planella, F. (2025). Comparing common zero-dimensional
SEI models under varied conditions. \emph{PyBaMM Conference 2025 Book
of Abstracts}, p.~24.

\bibitem[Mart\'{i}nez-Laserna et al.(2018)]{martinez2021second}
Mart\'{i}nez-Laserna, E., Gandiaga, I., Sarasketa-Zabala, E.,
Badeda, J., Stroe, D.-I., Swierczynski, M., \& Goikoetxea, A. (2018).
Battery second life: Hype, hope or reality? A critical review of the
state of the art. \emph{Renewable and Sustainable Energy Reviews},
\emph{93}, 701--718.

\bibitem[Nazir et al.(2025)]{nazir2025reservoirs}
Nazir, M. A., et al. (2025). PyBaMM-aided analysis of reservoirs in
lithium-ion cells: Strategies for cell design and cycling protocols.
\emph{PyBaMM Conference 2025 Book of Abstracts}, p.~25.

\bibitem[O'Kane et al.(2022)]{okane2022lithium}
O'Kane, S. E. J., Ai, W., Madabattula, G., Alonso-Alvarez, D., Timms, R.,
Sulzer, V., Edge, J. S., Wu, B., Offer, G. J., \& Marinescu, M. (2022).
Lithium-ion battery degradation: how to model it. \emph{Physical
Chemistry Chemical Physics}, \emph{24}(13), 7909--7922.

\bibitem[Ouledboutaarija et al.(2025)]{ouledboutaarija2025}
Ouledboutaarija, O., Kalogiannis, T., Hosen, M.~S., \& Berecibar, M.
(2025). Advancing second-life lithium-ion batteries: Optimization
techniques for improved electrochemical-ageing modeling performance.
\emph{PyBaMM Conference 2025 Book of Abstracts}, p.~19.

\bibitem[Peterson et al.(2010)]{peterson2010lithium}
Peterson, S. B., Apt, J., \& Whitacre, J. F. (2010). Lithium-ion
battery cell degradation resulting from realistic vehicle and
vehicle-to-grid utilization. \emph{Journal of Power Sources},
\emph{195}(8), 2385--2392.

\bibitem[Klass et al.(2019)]{klass2019variability}
Klass, V., Behm, M., \& Lindbergh, G. (2019). Capturing lithium-ion
battery dynamics with support vector machine-based battery model.
\emph{Journal of Power Sources}, \emph{298}, 92--101.

\bibitem[Planden et al.(2024)]{planden2024pybop}
Planden, B., Courtier, N.~E., Robinson, M., Khetarpal, A., Brosa
Planella, F., \& Howey, D.~A. (2024). PyBOP: A Python package for
battery model optimisation and parameterisation. \emph{arXiv preprint}
arXiv:2412.15859.

\bibitem[Prada et al.(2013)]{prada2013simplified}
Prada, E., Di Domenico, D., Creff, Y., Bernard, J., Sauvant-Moynot, V.,
\& Huet, F. (2013). Simplified electrochemical and thermal model of
LiFePO$_{4}$-graphite Li-ion batteries for fast charge applications.
\emph{Journal of The Electrochemical Society}, \emph{160}(4),
A616--A628.

\bibitem[Ramadass et al.(2004)]{ramadass2004firstprinciples}
Ramadass, P., Haran, B., Gomadam, P.~M., White, R., \& Popov, B.~N.
(2004). Development of first principles capacity fade model for Li-ion
cells. \emph{Journal of The Electrochemical Society}, \emph{151}(2),
A196--A203.

\bibitem[Sulzer et al.(2021)]{sulzer2021pybamm}
Sulzer, V., Marquis, S. G., Timms, R., Robinson, M., \& Chapman, S. J.
(2021). Python Battery Mathematical Modelling (PyBaMM). \emph{Journal
of Open Research Software}, \emph{9}(1), 14--21.

\bibitem[Weppner and Huggins(1977)]{weppner1977determination}
Weppner, W., \& Huggins, R. A. (1977). Determination of the kinetic
parameters of mixed-conducting electrodes and application to the system
Li$_{3}$Sb. \emph{Journal of The Electrochemical Society}, \emph{124}(10),
1569--1578.
\end{thebibliography}

\end{document}
